\section{Sard's Theorem}
\begin{theorem}[Sard's Theorem]\label{6.1}
    Let $f: U\to \mathbb{R}^n$ be a smooth map defined on an open subset $U\subset \mathbb{R}^m$. Then the set of critical values of $f$ ($f(C)$) has measure zero in $\mathbb{R}^n$.
\end{theorem}
\begin{remark}\label{6.2}
    \begin{itemize}
        \item If $m<n$, then $C=U$, so $f(U)$ has Lebesgue measure zero.
        \item $\mathbb{R}^n\setminus f(C)$ is dense in $\mathbb{R}^n$.
    \end{itemize}
\end{remark}
\begin{corollary}\label{6.3}
    Let $f: M\to N$ be a smooth map between smooth manifolds. Then the set of regular values of $f$ is dense in $N$.
\end{corollary}
\begin{proof}[Zard's Theorem]
    We proof the corollary by induction on $m$.\\
    If $m=0$, then $U\subseteq M$ is either empty or a single point, and the result is trivial.\\
    Now reminde yourself that $C$ is the set of critical points of $f$.\\
    We define $C_1:=\{x\in U|df_x=0\}$ and in similar fashion $C_i$ as the set of all $x\in U$ such that all partial derivatives of $f$ of order $\leq i$ vanish at $x$.\\
    It follows that $C\supseteq C_1\supseteq C_2\supseteq \ldots \supseteq C_i\supseteq C_{i+1}\supseteq \ldots$.\\
   
    Step 1: $f(C\setminus C_1)$ has measure zero.\\
        We may assume that $n\geq 2$ because for $n=1$, $C=C_1$.\\
        Now we take $\tilde x\in C\setminus C_1$. Our goal is to find an open neighbourhood $V\subseteq \mathbb{R}^n$ of $\tilde x$ such that $f(C\cap V)$ has measure zero.\\
        Since $\tilde x\notin C_1$, there exists a partial derivative of $f$, say $\frac{\partial f_1}{\partial x_1}$, that does not vanish at $\tilde x$.\\
        Now define $h:U\to \mathbb{R}^n$ by $h(x)=(f_1(x),x_2,\ldots,x_m)$.\\
        This is a local diffeomorphism at $\tilde x$ because at $\tilde x$ $dh_{\tilde x}$ is non-singular, hence it maps a neighbourhood $V$ of $\tilde x$ diffeomorphically onto an open set $V'\subseteq \mathbb{R}^m$.\\
        We also define $g:V'\to \mathbb{R}^n$ by $g(y)=f\circ h^{-1}$.\\
        Let $C'$ be the set of critical points of $g$. Then $C'=h(C\cap V)\implies g(C')=(g\circ h)(V\cap C)=f(V\cap C)$.\\
        For each $(t,x_2,...,x_m)\in V'$, we have $g(t,x_2,...,x_m)\in \{t\}\times \mathbb{R}^{n-1}$, hence $g$ maps a hyperplane to a hyperplane.\\
        Now we define $g^t:(\{t\}\times \mathbb{R}^{m-1})\cap V'\to \{t\}\times \mathbb{R}^{n-1}$ the restriction of $g$.
        The Jacobian matrix of $g^t$ is $(\frac{\partial g_i}{\partial x_j})=\begin{cases}
            1 & 0\\
            * & \frac{\partial g_i}{\partial x_j}
        \end{cases}$ 
        Hence a point in $\{t\}\times \mathbb{R}^{n-1}$ is critical for $g$ iff it is critical for $g^t$.\\
        By induction the set of critical values of $g^t$ has measure zero in $\{t\}\times \mathbb{R}^{n-1}$.\\
        With Fubini's theorem, it follows that $g(C')=f(V\cap C)$ has measure zero in $\mathbb{R}^n$.\\
    Step 2: $f(C_i\setminus C_{i+1})$ has measure zero.\\      
    
\end{proof}
